%==============================================================================
% الخلاصة (Summary) — official project abstract per Ministry of Higher
% Education and Scientific Research guide. Text approved by the team.
%==============================================================================
\chapter*{الخلاصة}
\addcontentsline{toc}{chapter}{الخلاصة}
\markboth{الخلاصة}{الخلاصة}

تواجه المؤسسات الأكاديمية والإدارية في البيئة العربية عموماً، واليمنية خصوصاً،
تحديات كبيرة في الوصول السريع والذكي إلى بياناتها التشغيلية ومستنداتها
التنظيمية. ويؤدي الاعتماد على الأساليب اليدوية التقليدية والافتقار إلى أنظمة
ذكية تدعم اللغة العربية إلى هدر الموارد، وبطء في اتخاذ القرار، وصعوبة في
استخراج التقارير والرؤى الاستراتيجية.

لمعالجة هذه الفجوة البحثية والتطبيقية، يقدم هذا المشروع «بيان»، وهو مساعد
ذكي مؤسسي متكامل يعتمد على معمارية الأنظمة متعددة الوكلاء
(\EN{Multi-Agent System}) لتوفير واجهة محادثة تفاعلية باللغة العربية (الفصحى
والعامية). يدمج النظام بشكل مبتكر بين أربع تقنيات متقدمة في إطار موحد: تحويل
اللغة الطبيعية إلى استعلامات قواعد بيانات (\EN{NL2SQL})، الاسترجاع المعزز
بالتوليد (\EN{RAG}) للبحث الدلالي في المستندات، معالجة اللغة العربية
(\EN{Arabic NLP})، وتنسيق الوكلاء البرمجيين عبر إطار \EN{LangGraph}.

وللتغلب على تحديات تعقيد قواعد البيانات وظاهرة «الهلوسة» في النماذج اللغوية،
يقدم المشروع مساهمة ابتكارية تتمثل في تصميم «طبقة الرؤى الذكية»
(\EN{Smart Views})، وهي طبقة تجريدية تبسط مخطط قاعدة البيانات لرفع دقة توليد
الاستعلامات. تم توجيه وتصميم النظام ليعمل بالتكامل مع بيئة
«\EN{Odoo Education}» كبيئة تطبيقية واقعية ومحاكاة للأنظمة الجامعية.

يهدف النظام إلى تمكين الموظفين وصناع القرار من الاستعلام اللحظي عن البيانات
الآمنة (المهيكلة)، والبحث الموثق في اللوائح والسياسات (غير المهيكلة)،
واستخراج المؤشرات الإحصائية وتوليد التقارير بصورة مؤتمتة تماماً. وتتوقع
الدراسة أن يسهم هذا النظام في تسريع العمليات التشغيلية، وتعزيز الشفافية
والموثوقية عبر سجلات التدقيق (\EN{Audit Logs}) وتتبع مسار وكلاء الذكاء
الاصطناعي (\EN{Agent Trace})، مما يقدم نموذجاً فعالاً وقابلاً للتوسع يدعم
جهود التحول الرقمي في المؤسسات المحلية.

\clearpage
